\section{Model, training and the feed-forward control}
\label{app:model}

\paragraph{Architecture.}


Table~\ref{tab:arch} lists the network layer by layer. The tile crop is one-hot
over nine tile types with two coordinate channels appended, the CoordConv
construction of \citet{liu2018intriguing}; three valid $3\times3$
convolutions reduce it to $15\times15\times32$, the flattened features and the
five scalars are projected to a $256$-d embedding, and a single-layer GRU with a
$128$-d state carries the only memory. Two linear heads read the state. There is
no other head and no other loss: the map category is never given to the network
as a target, so any belief in $\mathbf{h}_t$ is learned from reward alone.

\begin{figure}[h]
    \centering
    % PPO+GRU agent. Included from sections/method.tex.
\resizebox{\linewidth}{!}{%
\begin{tikzpicture}[
  font=\scriptsize,
  box/.style={draw, rounded corners=2pt, fill=gray!8, minimum height=10mm,
              inner sep=3pt, align=center},
  arr/.style={-{Latex[length=2mm]}, thick},
  node distance=5mm and 5mm]
  \node[box] (obs) {egocentric\\view};
  \node[box, right=of obs] (enc) {one-hot\\$+$ CoordConv};
  \node[box, right=of enc] (cnn) {$3\times$ convolution};
  \node[circle, draw, inner sep=1pt, right=6mm of cnn] (cat) {$\oplus$};
  \node[box, right=6mm of cat] (mlp) {Linear ReLU\\$7205\to256$};
  \node[box, right=of mlp, fill=gray!25, thick] (gru) {GRU\\$\mathbf{h}_t\in\mathbb{R}^{128}$};
  \node[below=5mm of cat] (sc) {scalars};
  \coordinate[right=9mm of gru] (split);
  \node[box, above right=2.5mm and 7mm of split] (actor) {Linear Actor\\$\pi$};
  \node[box, below right=2.5mm and 7mm of split] (critic) {Linear Critic\\$V$};
  \node[right=3mm of actor] (acts) {$6$ actions};
  \draw[arr] (obs) -- (enc);
  \draw[arr] (enc) -- (cnn);
  \draw[arr] (cnn) -- (cat);
  \draw[arr] (cat) -- (mlp);
  \draw[arr] (sc) -- (cat);
  \draw[arr] (mlp) -- (gru);
  \draw[thick] (gru) -- (split) node[midway, above] {$\mathbf{h}_t$};
  \draw[arr] (split) |- (actor);
  \draw[arr] (split) |- (critic);
  \draw[arr] (actor) -- (acts);
  \fill (split) circle (1pt);
  \draw[arr, dashed] (split) ++(-4.5mm,0) -- ++(0,-10mm) -| (gru.south)
        node[pos=0.25, below] {$\mathbf{h}_{t-1}$};
\end{tikzpicture}}
\caption{\textbf{The agent.} The $21\times21$ tile view is one-hot encoded with
    two coordinate channels appended \citep{liu2018intriguing}, passed through
    three layers of convolutions, concatenated with five scalars (the facing direction as a one-hot of four and the elapsed time).
    This vector is projected to a
    $256$-d embedding, which is fed to a GRU with $128$-d hidden state $\mathbf{h}_t$. The actor and critic are linear projections of $\mathbf{h}_t$.}
    \label{fig:model}
\end{figure}
\begin{table}[h]
\centering
\small
\caption{\textbf{Architecture of the PPO+GRU agent.} The two coordinate
channels follow \citet{liu2018intriguing}. Weights are initialised
orthogonally, the actor with gain $0.01$. $2{,}015{,}559$ parameters in total.}
\label{tab:arch}
\begin{tabular}{@{}llrr@{}}
\toprule
Layer & Detail & Output & Parameters \\
\midrule
Input & $21\times21$ crop, one-hot $9$ tiles $+\,2$ CoordConv channels & $21\times21\times11$ & -- \\
Conv 1 & $3\times3$, valid, ReLU & $19\times19\times32$ & $3{,}200$ \\
Conv 2 & $3\times3$, valid, ReLU & $17\times17\times32$ & $9{,}248$ \\
Conv 3 & $3\times3$, valid, ReLU & $15\times15\times32$ & $9{,}248$ \\
Embedding & flatten $(7{,}200)$ $\oplus$ $5$ scalars $\to$ linear, ReLU & $256$ & $1{,}844{,}736$ \\
GRU & single layer, state $\mathbf{h}_t$ & $128$ & $148{,}224$ \\
Actor & linear, categorical over $6$ actions & $6$ & $774$ \\
Critic & linear & $1$ & $129$ \\
\bottomrule
\end{tabular}
\end{table}

\paragraph{Training.}
Table~\ref{tab:hparams} gives every hyperparameter. The loss is the clipped
PPO objective with a value term and an entropy bonus \citep{schulman2017ppo},
advantages from generalised advantage estimation \citep{schulman2016gae}, and
backpropagation through time over each $128$-step rollout. The entropy
coefficient starts at $0.15$ and is annealed linearly to zero. This schedule is
what makes the plain reward solvable: with a small entropy bonus the policy
settles on a constant flag, which already pays on two thirds of maps, and
never learns to use the category. The released agent is the best of five
belief-free seeds at entropy $0.15$; four solved the task and one collapsed
to the constant flag (Figure~\ref{fig:training}, Table~\ref{tab:seeds}). Runs at entropy $0.10$ and $0.20$ also solved the
task, so the recipe is not fragile in that knob. An otherwise equal agent
trained with an auxiliary category head scores the same held out, so
supervising the belief buys nothing.

\begin{figure}[h]
\centering
\includegraphics[width=\linewidth]{figures/training_curves.png}
\caption{\textbf{Training curves, mean $\pm$ one standard deviation over five
seeds.} Return per episode, success on all maps, and success on the decisive
lakes and rocky maps, against environment steps. The feed-forward control
plateaus at the constant-flag optimum: two thirds of maps overall and chance
on the decisive ones. The recurrent agent's band is wide early and late
because one of its five seeds collapses to the same constant flag; the others
reach near-perfect success.}
\label{fig:training}
\end{figure}

\begin{table}[h]
\centering
\small
\caption{\textbf{Training hyperparameters.} Identical for the recurrent agent
and the feed-forward control.}
\label{tab:hparams}
\begin{tabular}{@{}lr@{}}
\toprule
Hyperparameter & Value \\
\midrule
Environment steps & $6\times10^6$ \\
Parallel environments $\times$ rollout length & $32\times128$ \\
Minibatches per update $\times$ epochs & $4\times4$ \\
Optimiser, learning rate & Adam, $3\times10^{-4}$, constant \\
Discount $\gamma$, GAE $\lambda$ & $0.99$, $0.95$ \\
Clip $\epsilon$ & $0.2$ \\
Value coefficient & $0.5$ \\
Entropy coefficient & $0.15 \to 0$, linear over training \\
Gradient-norm clip & $0.5$ \\
Observation encoding & one-hot tiles, $21\times21$ view \\
Map pool & $4{,}800$ training maps, categories sampled uniformly \\
Evaluation & sampled actions, $1{,}200$ held-out maps \\
\bottomrule
\end{tabular}
\end{table}

\paragraph{Compute.}
Each run used one NVIDIA RTX~3090 of a shared SLURM cluster, with the $32$
environments stepped on the CPU, at about $2{,}450$ environment steps per
second for the recurrent agent and $8{,}000$ for the feed-forward control:
roughly $40$ minutes for a recurrent run and $13$ minutes for a feed-forward
run of $6$\,M steps. The whole sweep of twelve runs fits in a few GPU-hours.

\paragraph{Feed-forward control.}
To separate navigation from memory we train the same network with the GRU
replaced by a per-step linear layer with ReLU of the same width, so the agent
sees exactly the same input and has exactly the same heads but no carried
state. Five seeds with the same schedule as the recurrent agent, $6$\,M steps
each; all converge within the first million steps (Figure~\ref{fig:training}). Table~\ref{tab:seeds} gives the held-out result
of every seed of both agents. The feed-forward agent learns navigation: it
reaches a flag on every map and crosses water and rock as needed. What it
cannot do is choose the flag: each seed settles on one constant flag, so it is
right on all balanced maps and on one of the two decisive categories, at
chance on the decisive maps as a whole. The gap between the two agents is
therefore the memory itself, the belief integrated across the corridor.

\begin{table}[h]
\centering
\small
\caption{\textbf{Held-out results of every seed}, $1{,}200$ test maps, one
sampled rollout per map; both agents trained for $6$\,M steps with identical
hyperparameters. \emph{Decisive} is success on lakes and rocky maps,
where chance is $50\%$; \emph{top} is the share of balanced maps exited by the
top flag. Every feed-forward seed reaches a flag on essentially every map but
always the same one, so it scores $100\%$ on balanced maps, $100\%$ on one
decisive category and $0\%$ on the other. Four of five recurrent seeds solve
the task; seed~1 fell into the same constant-flag optimum. $^\dagger$The
released agent. Evaluations sample actions, so repeated runs of one checkpoint
differ by a few tenths of a point.}
\label{tab:seeds}
\begin{tabular}{@{}lrrrrrr@{}}
\toprule
& Overall & Decisive & Balanced & Lakes & Rocky & Top (balanced) \\
\midrule
\multicolumn{7}{@{}l}{\emph{GRU (recurrent)}} \\
seed 1 & 66.4 & 49.8 & 99.8 & 99.5 & 0.0 & 0.0 \\
seed 2$^\dagger$ & 98.8 & 98.2 & 100.0 & 98.5 & 98.0 & 57.2 \\
seed 3 & 97.8 & 97.0 & 99.2 & 94.2 & 99.8 & 64.0 \\
seed 4 & 98.5 & 97.8 & 100.0 & 97.5 & 98.0 & 41.2 \\
seed 5 & 97.2 & 96.1 & 99.5 & 98.2 & 94.0 & 71.0 \\
mean $\pm$ std & 91.8 $\pm$ 12.7 & 87.8 $\pm$ 19.0 & 99.7 $\pm$ 0.3 & 97.6 $\pm$ 1.8 & 78.0 $\pm$ 39.0 & 46.7 $\pm$ 25.3 \\
\addlinespace
\multicolumn{7}{@{}l}{\emph{Feed-forward control}} \\
seed 1 & 66.5 & 49.9 & 99.8 & 99.8 & 0.0 & 0.0 \\
seed 2 & 66.5 & 49.9 & 99.8 & 0.0 & 99.8 & 99.8 \\
seed 3 & 66.2 & 49.9 & 99.0 & 0.0 & 99.8 & 99.0 \\
seed 4 & 66.5 & 49.9 & 99.8 & 0.0 & 99.8 & 99.8 \\
seed 5 & 66.4 & 49.8 & 99.8 & 0.0 & 99.5 & 99.8 \\
mean $\pm$ std & 66.4 $\pm$ 0.1 & 49.9 $\pm$ 0.1 & 99.6 $\pm$ 0.3 & 19.9 $\pm$ 39.9 & 79.8 $\pm$ 39.9 & 79.7 $\pm$ 39.8 \\
\bottomrule
\end{tabular}
\end{table}



